\Chapter{Stopień realizacji celu i wnioski}
\label{ch:conclusions}

Cel końcowy zrealizowano częściowo: przygotowano i przetestowano
fundament danych, lecz nie zbudowano jeszcze systemu przewidującego
wymianę opon. Wykonane są kontrakty, konfiguracja środowiska, testy
offline, audyt FastF1, rzeczywiste fixtures oraz eksporter i weryfikator
snapshotu. Kod T2 jest włączony do \code{main}; akceptacja danych
i publikacja archiwum są osobnymi, oczekującymi działaniami.

Audyt przyniósł konkretny wniosek metodologiczny: walidację znaczenia
odpowiedzi należy wykonać przed treningiem. W tym projekcie wzrost
numeru stintu okazał się niewystarczającym dowodem wymiany opon,
a końcowa liczba planowanych okrążeń wymaga interpretacji czasowej.
Zachowanie tych niepewności w danych i raporcie jest bardziej użyteczne
niż przypisanie modelowi pozornej precyzji na błędnym celu.

Najbliższa praca obejmuje formalne domknięcie pochodzenia źródeł w T2c
i akceptację snapshotu. Przed etykietowaniem potrzebna jest decyzja D2,
a przed odpowiednią cechą także kontrakt D5 i polityka sprzed pierwszej
zapowiedzi liczby okrążeń. Następnie można zaimplementować etykiety,
podziały, przyczynowe cechy, standaryzację i okna wraz z testami wycieku.
Dopiero ten zbiór stanowi podstawę uczciwego porównania modeli i panelu.

Nie ma jeszcze wyników, które pozwalałyby ocenić skuteczność predykcji.
Końcowe wnioski o metodach i zastosowaniach będą wymagały artefaktów
uczenia, jawnej oceny walidacyjnej oraz osobno autoryzowanej oceny
testowej. Bieżący dokument stanowi merytoryczny opis ukończonej części
projektu i punkt wyjścia do aktualizacji po przyjęciu kolejnych etapów.
